% a04_Bh_bulk3  Proof of app:a04:lem:B2, part 3: row (III) with rate, the sup bound, the original-vertex part assembled

\subsubsection*{B2.6 Row $(\mathrm{III})$ with rate, the sup bound, and the original-vertex part}

\begin{theorem}[row $(\mathrm{III})$]\label{app:a04:thm:III}
Let $\cT_\tau(b)\defeq\tau^3\sum_{x\in\Lambda_\tau\cap[0,1)^3}\big(G(\Pi_C-\Pi_G)G\big)_\tau(x,x)$ (all lines at spacing $\tau$ with block coupling $b$ and the weight of that member). For every $\theta<1$ there is $C_\theta$, depending on $\theta$, $\mathfrak a$, $m_0$ and $\norm{f}_\infty$ only, such that for all $t=5^{-j}\le1$ ($\le\frac15$ for $f_D$) and $m',\mu'\in(0,m_0]$:
\[
\textup{(III-sup)}\ \ \sup_x\big|\big(G(\Pi_C-\Pi_G)G\big)_t(x,x)\big|\le C(1+\ln t^{-1})^{14};\qquad
\textup{(III-inc)}\ \ \big|\cT_{t/5}(\mathfrak a_{j+1})-\cT_t(\mathfrak a_j)\big|\le C_\theta t^\theta.
\]
\end{theorem}
\begin{proof}
By \cref{app:a04:lem:sandwich,app:a04:lem:D}, $(G(\Pi_C-\Pi_G)G)(x,x)$ is a finite sum (independent of $\tau$) of triangle terms $T_\tau(x)=\tau^6\sum_{u,v}k_1(x,u)k_2(u,v)k_3(v,x)$ and of two-point (seagull) terms, with $k_1,k_3\in\{G_\tau,\nabla G_\tau\}$ (slots $a,c\in\{1,2\}$), $k_2\in\{\Gamma D,\Gamma\nabla D,\nabla\Gamma D,G_{\mathrm{in}}\Gamma^\ph,\nabla G_{\mathrm{in}}\Gamma^\ph,G_{\mathrm{in}}\nabla\Gamma^\ph\}$ (slot $b\in\{1,2\}$), $a+b+c\le5$, arguments shifted by at most one lattice unit. Put $R\defeq R_{t/5}+5\le C(1+\ln t^{-1})$ and split each outer line $k_i=k_i^{\mathrm{in}}+k_i^{\mathrm{out}}$, $k^{\mathrm{in}}\defeq k\mathbf 1_{|r|\le R}$. \emph{Sup bounds} (L1): $|k_1|\le C(R+1)^aw_a$, $|k_3|\le C(R+1)^cw_c$ on $|r|\le R$ and $\le C\tau^4e^{-c|r|}$ on $|r|\ge R-5$ (O1), (O2); $|k_2|\le Cw_b$ (($\Gamma$1), ($\Gamma$2), (T1), (O1)), with $C$ replaced by $C(2R+1)^b$ for an inner $G_{\mathrm{in}}$, $\nabla G_{\mathrm{in}}$ on the confined region $|u-v|\le2R$ (a polylogarithm, absorbed below); where an outer line is ``out'', such an inner line is bounded by $Cw_b(u-v)$ on $|u-v|\le1$ and by $C$ beyond (slot $0$). The weight enters only through \cref{app:a04:lem:Gam,app:a04:lem:TIOUT}; the kernels are evaluated, never their weights.

\emph{Step 0 (sup bound).} Core (both outer lines ``in''): \cref{app:a04:lem:PC} with $\Sigma\le5$ gives $\le C(R+1)^{a+c}(2R+1)^b(1+\ln(R/\tau))^4R^{6-\Sigma}\le C(1+\ln t^{-1})^{13}$. Terms with $k_1^{\mathrm{out}}$: $\le C\tau^4\sum_u\tau^3e^{-c|x-u|}[\tau^3\sum_vw_b(u-v)|k_3(v,x)|]$; by \cref{app:a04:lem:Tp} (the $v$-sum over the ball $|v-x|\le R$) and the $k_3^{\mathrm{out}}$ tail, the bracket is $\le C(R+1)^{c+2}(1+\ln\tau^{-1})(1+w_{(b+c-3)_+}(u-x))+C\tau^4(1+\ln\tau^{-1})$ (the case $b+c<3$ of \cref{app:a04:lem:Tp}, with slot $0$ for an inner $G_{\mathrm{in}}$ beyond $|u-v|\le1$, gives $R^{3-b-c}\le R^2$), and the tail version of \cref{app:a04:lem:PC} gives $\le C\tau^4(1+\ln\tau^{-1})^{q}$. Same for $k_3^{\mathrm{out}}$. Seagull terms: the two-point version of \cref{app:a04:lem:PC}. Hence (III-sup) with exponent $14$, uniformly in $m',\mu'$: no line needs mass decay, the outer lines decaying at the $m'$-independent rate $c$.

\emph{Step 1 (decomposition of the increment).} With the cells $c(x)\defeq x+(t/5)\{0,\dots,4\}^3$ (they tile $\Lambda_{t/5}$, and the fine unit cell is the disjoint union of the $c(x)$, $x\in\Lambda_t\cap[0,1)^3$), for each triangle term
\[
\cT_{t/5}(\mathfrak a_j)-\cT_t(\mathfrak a_j)=\cD+\cC,\quad
\cD\defeq t^9\!\!\sum_{x,u,v\in\Lambda_t}\!\Big[5^{-9}\!\!\sum_{x'\in c(x),u'\in c(u),v'\in c(v)}\!\!F_{t/5}(x',u',v')-F_{t/5}(x,u,v)\Big],\quad
\cC\defeq t^9\!\!\sum_{x,u,v\in\Lambda_t}\![F_{t/5}-F_t](x,u,v),
\]
$F_\tau(x,u,v)\defeq k_1k_2k_3$ at spacing $\tau$ ($x$ restricted to the unit cell). Near-diagonal set $\cN\defeq\{\min(|x-u|,|u-v|,|v-x|)<8t\}$. (For $f_D$ the fine member's weight lives on the $t$-blocks, which are exactly the cells $c(x)$, the coarse member's on the $5t$-blocks.)

\emph{Step 2 ($\cN$).} Each of the four sums restricted to $\cN$ (fine triples whose coarse parent is in $\cN$ lie in $\cN'\defeq\{\min<8t+2\sqrt3t\}$) is bounded by (L1) alone. Case $|x-u|<8t$: $\le17^3$ values of $u$ (resp.\ $(17\cdot5)^3$ fine ones), $|F|\le C(R+1)^{a+c}\tau^{-a}w_b(u-v)w_c(v-x)$; \cref{app:a04:lem:Tp} ($R$-confinement by $k_3$, or its tail) gives $\tau^3\sum_v\le C(1+\ln\tau^{-1})R^3t^{-(b+c-3)_+}$, so the contribution is $\le Ct^3t^{-a}t^{-(b+c-3)_+}(1+\ln t^{-1})^{8}\le Ct(1+\ln t^{-1})^{8}$, since $a\le2$ and $a+b+c\le5$ give $3-a-(b+c-3)_+\ge\min(3-a,6-\Sigma)\ge1$. Case $|v-x|<8t$: symmetric. Case $|u-v|<8t$: $\le17^3$ values of $v$ per $u$, $|F|\le C(R+1)^{a+c}w_a(x-u)t^{-b}w_c(v-x)$ with $w_c(v-x)\le2^cw_c(u-x)$ if $|u-x|\ge16t$ and $\le t^{-c}$ otherwise: $\le Ct^{3-b}(R+1)^{a+c}[t^3\sum_{|x-u|\le R}w_{a+c}(x-u)+t^3\cdot33^3t^{-a-c}]\le Ct(1+\ln t^{-1})^{8}$ ($b\le2$, $a+c\le4$, $\Sigma\le5$). Two-point terms: the same with one fewer sum.

\emph{Step 3 ($\cC$ off $\cN$: same points, two spacings).} $F_{t/5}-F_t=\sum_\ell(k_\ell^{t/5}-k_\ell^t)\prod_{\text{others}}$ (mixed products, every factor bounded by (L1) at its own spacing). Line differences at common points $x,u,v\in\Lambda_t$ with all distances $\ge8t$: translation-invariant factors $C,\nabla C,D,\nabla D$: $\le Ctw_{s+1}$ by (T2), slot $s\to s+1$; coarse factors $\Gamma,\Gamma^\ph$: $\le Ct$, $\nabla\Gamma,\nabla\Gamma^\ph$: $\le Ct(1+\ln t^{-1})^3$ by ($\Gamma$4), slots unchanged (for $f_D$, each member with its own weight, ($\Gamma$4) as stated); outer lines $G=C-\Gamma$, $\nabla G$: the sum of a translation-invariant and a coarse difference, hence $\le Ctw_{s+1}+Ct(1+\ln t^{-1})^3\le Ct(1+\ln t^{-1})^3(R+1)^{s+1}w_{s+1}$ on $|r|\le R$, and $\le2C\tau^4e^{-c|r|}$ on $|r|\ge R$ (both members are tails there, (O2), with the same $\alpha_0$ for both spacings). An argument shifted by one lattice unit at the two spacings ($te_\mu$ versus $(t/5)e_\mu$) is a displacement $\le t$ and falls under Step~4. For a differenced inner line the two outer lines still confine $u,v$ to $R$-balls (or are tails); for a differenced outer line its ``in'' part is confined by definition and its ``out'' part carries $e^{-c|r|}$: no mass decay is used. \cref{app:a04:lem:PC} with $\Sigma\le6$, $R\le C(1+\ln t^{-1})$, and the tail version: $|\cC\text{ off }\cN|\le Ct(1+\ln t^{-1})^{q}$.

\emph{Step 4 ($\cD$ off $\cN$: displaced arguments, one spacing $\tau=t/5$).} For $(x',u',v')\in c(x)\times c(u)\times c(v)$ every pairwise displacement is $\le2\sqrt3t\le|r|/4$ ($|r|\ge8t$), so $|r'|\ge|r|/2$ and (L1) at the displaced arguments is $\le2^s$(L1) at $(x,u,v)$; $R$-confinement holds with $R\pm4$. Telescope over the lines: translation-invariant factors by (T3): $\le Ctw_{s+1}(r)$; coarse factors $\Gamma,\Gamma^\ph$ by ($\Gamma$2) along lattice paths: $\le Ct$; $\nabla\Gamma,\nabla\Gamma^\ph$ by ($\Gamma$3): $\le Ct^\theta$ (slots unchanged; for $f_D$ the finest-scale modulation of $\nabla\Gamma$ is inside the constant of ($\Gamma$3)); outer lines: translation-invariant plus coarse parts as above, tails by (O2). \cref{app:a04:lem:PC} ($\Sigma\le6$): $|\cD\text{ off }\cN|\le C_\theta t^\theta(1+\ln t^{-1})^{q}$. Two-point terms: the same.

\emph{Step 5 ($b$).} $|\cT_{t/5}(\mathfrak a_{j+1})-\cT_{t/5}(\mathfrak a_j)|$: telescope over the lines; translation-invariant lines do not depend on $b$. Inner coarse factors ($\Gamma,\Gamma^\ph$, and $G_{\mathrm{in}}$ through $\Gamma$): ($\Gamma$5) gives $\le C|\mathfrak a_j-\mathfrak a_{j+1}|\le Ct^2$, confined by the two outer lines as in Steps~3--4. A differenced outer line: the same in/out dichotomy. \cref{app:a04:lem:CT} is uniform in $b\in[\mathfrak a(1-5^{-2}),\mathfrak a]$, so for $|r|\ge R$ both members obey (O2), $|k(\mathfrak a_{j+1})(r)-k(\mathfrak a_j)(r)|\le2C\tau^4e^{-c|r|}$; for $|r|\le R$ use ($\Gamma$5): $\le Ct^2(R+1)^sw_s(r)$, $s$ its slot. \cref{app:a04:lem:PC} and its tail version give $\le Ct^2(1+\ln t^{-1})^q$.

Summing Steps~2--5 and absorbing the logarithms into $t^\theta$ ($\theta<1$ arbitrary) gives (III-inc). Summation by parts (\cref{app:a04:lem:D}) was on the infinite lattice, so there is no boundary term, and all sums are absolutely convergent at fixed $\tau$ (outer lines decay exponentially by (O2), inner lines are bounded by $w_1$ at worst).
\end{proof}

\begin{proposition}[the sup bound at order $e^2$]\label{app:a04:prop:S}
For both weights, for all $k$ and $1\le j\le k$, with $C$ independent of $k$ and $j$ (depending on $m$, $\mathfrak a$, $m_0$, $\norm f_\infty$):
\[
\big\|e^2[\mathfrak g_x(C,C,C;D)-\mathfrak g_x(G_j,G_j,G_j;D_j)]\big\|_\infty=\norm{\beta_{(2,0)}+\mathfrak a}_\infty\le Ce^2\big[(1+(m\eta)^{-1})+(1+j\ln5)^{14}\big]\le Ce^2\eta^{-1}(1+j)^{14}.
\]
\end{proposition}
\begin{proof}
Pointwise (no cell average), at unit block scale with $\tau=t=5^{-j}$: $(\Gamma\Pi_CC)(x,x)=\int_{\mathbb T}\Mh\,\overline{\hat v_x}\,\hat u_x\dP$ with $\hat u_x\defeq\sum_{n\in N_\tau}s_\tau\hat C_\tau^2\Pih_\tau(p)e^{-ip\cdot x}$; by (E1)--(E3) and ($\Pi$-a), $|\hat u_x|\le C(|P|+m_0)(P^2+m'^2)^{-2}+C\sum_{n\ne0}\rho(n)|n|^{-3}\le C(|P|+m_0)(P^2+m'^2)^{-2}$ (for $f_D$ the D-aliases add $\sum_{\ell\ne0}\sigma(q)(25\tau^2/4)^2C\sqrt3\pi/\tau\le C\tau^3(1+\ln\tau^{-1})^3$), and $|\hat v_x|\Mh\le C$ ((A3) for $f_D$); so $|(\Gamma\Pi_CC)(x,x)|\le C\int_{\mathbb T}(|P|+m_0)(P^2+m'^2)^{-2}\dP\le C(1+\ln(1/m')+m_0/m')\le C(1+m'^{-1})$ (this is the infrared logarithm of the order-$e^2$ Wick constant). $(C\Pi_C\Gamma)(x,x)$: the same. $(\Gamma\Pi_C\Gamma)(x,x)=\int\Mh^2|\hat v_x|^2Y_\tau\dP\le C\int(P^2+m'^2)^2(P^2+m'^2)^{-2}(|P|+m_0)(P^2+m'^2)^{-2}\dP\le C(1+m'^{-1})$. $(\mathrm{III})$: (III-sup). By \eqref{app:a04:eq:dens}--\eqref{app:a04:eq:sandwich} the original-vertex density is $-2e^2[(\mathrm I)+(\mathrm{II})+(\mathrm{III})]$ at unit block scale; with $m'=m\eta$, $1+m'^{-1}=1+(m\eta)^{-1}$ and $(1+\ln t^{-1})^{14}=(1+j\ln5)^{14}$. The last inequality uses $\eta\le1$ and $(m\eta)^{-1}\le m^{-1}\eta^{-1}$.
\end{proof}

\begin{theorem}[the original-vertex part]\label{app:a04:thm:B2orig}
Let $\theta\in(0,1)$. There are $C$, $C_\theta$ (depending on $\theta$, $\mathfrak a$, $m_0\ge\max(m,\lmu)\eta$, $e$, $m$ and $\norm f_\infty$) such that for all $k$ and $1\le j\le k$ ($t=5^{-j}$, $\eta=5^j\eps_k\le1$), for both weights,
\begin{align*}
\norm{\beta_{(2,0)}+\mathfrak a}_\infty&\le Ce^2\big[(1+(m\eta)^{-1})+(1+j\ln5)^{14}\big]\le Ce^2\eta^{-1}(1+j)^{14},\\
\big|(\bar\beta_{(2,0)}+\bar{\mathfrak a})^{(j)}_k-(\bar\beta_{(2,0)}+\bar{\mathfrak a})^{(j+1)}_{k+1}\big|&\le C_\theta e^2\big[t(1+(m\eta)^{-1})+t^\theta\big]\le C_\theta e^2\,5^{-\theta j}\eta^{-1},
\end{align*}
where $\bar{\mathfrak a}$ is the block mean of $\mathfrak a_x$. This is the statement of \cref{app:a04:lem:B2} for the original-vertex part, with $\vartheta=\theta$ arbitrary in $(0,1)$ and $q=14$.
\end{theorem}
\begin{proof}
By \eqref{app:a04:eq:dens}--\eqref{app:a04:eq:sandwich}, $\beta_{(2,0)}+\mathfrak a=-2e^2[(\mathrm I)+(\mathrm{II})+(\mathrm{III})]$ at unit block scale, and the block mean $\eta^{-3}\sum_{x\in\Lk\cap[0,\eta)^3}\eps_k^3(\cdot)$ is the cell average $\langle\cdot\rangle_\cell$ of \cref{app:a04:lem:cell} at unit block scale. The diagonal pair $(k,j)\to(k+1,j+1)$ keeps $\eta,m',\mu'$ and the unit cell fixed and changes $(\tau,b,\text{weight})=(t,\mathfrak a_j,f\text{ on }5t\text{-blocks})\to(t/5,\mathfrak a_{j+1},f\text{ on }t\text{-blocks})$ (the reduction of \cref{app:a04:lem:B0} and its Dirichlet-weight paragraph). \Cref{app:a04:thm:R}: $|[(\mathrm I)+(\mathrm{II})]_{t/5}(\mathfrak a_{j+1})-[(\mathrm I)+(\mathrm{II})]_t(\mathfrak a_j)|\le Ct(1+m'^{-1})$. \Cref{app:a04:thm:III}: $|\cT_{t/5}(\mathfrak a_{j+1})-\cT_t(\mathfrak a_j)|\le C_\theta t^\theta$. Sup bound: \cref{app:a04:prop:S}. With $m'=m\eta$: $t(1+m'^{-1})\le(1+m^{-1})5^{-j}\eta^{-1}$ and $t^\theta=5^{-\theta j}\le5^{-\theta j}\eta^{-1}$. For $f_D$ the restriction $j\ge1$ ($\tau\le\frac15$ for both members) is the one in the statement of \cref{app:a04:lem:B2}.
\end{proof}
